% Part: incompleteness
% Chapter: representability-in-q
% Section: representing-relations

\documentclass[../../../include/open-logic-section]{subfiles}

\begin{document}

\olfileid{inc}{req}{rel}

\olsection{د اړيکو تمثيلول}

دا بيان کړو چې يوه \emph{اړيکه} تمثيلېدونکې وي څه معنا لري۔

\begin{defn}
\ollabel{defn:representing-relations} پر طبيعي عددونو اړيکه
$R(x_0,\dots,x_k)$ هله \emph{په~$\Th{Q}$ کښې تمثيلېدونکې} ده چې
داسې فارمول~$!A_R(x_0,\dots,x_k)$ وي چې هر کله
$R(n_0,\dots,n_k)$ رښتيا وي، $\Th{Q}$،
$!A_R(\num{n_0},\dots,\num{n_k})$ ثابت کړي، او هر کله
$R(n_0,\dots,n_k)$ ناسم وي، $\Th{Q}$،
$\lnot !A_R(\num{n_0}, \dots, \num{n_k})$ ثابت کړي۔
\end{defn}

\begin{thm}
\ollabel{thm:representing-rels} يوه اړيکه په~$\Th{Q}$ کښې هله او
يوازې هله تمثيلېدونکې ده چې محاسبه کېدونکې وي۔
\end{thm}

\begin{proof}
د مخکښ لوري دپاره فرض کړئ اړيکه~$R(x_0,\dots,x_k)$ په
فارمول~$!A_R(x_0,\dots,x_k)$ تمثيلېږي۔ د~$R$ د محاسبې الګوريتم دا
دے: پر ننوتونو~$n_0$، \dots،~$n_k$ په هممهاله ډول د
$!A_R(\num{n_0}, \dots, \num{n_k})$ د ثبوت او د
$\lnot !A_R(\num{n_0}, \dots, \num{n_k})$ د ثبوت پلټنه وکړئ۔ زموږ
د فرض له مخې پلټنه خامخا يو له دغو دواړو مومي؛ که لومړے وي، «هو»
راپور کړئ، او که نۀ، «نه» راپور کړئ۔

په بل لوري کښې فرض کړئ~$R(x_0, \dots, x_k)$ محاسبه کېدونکې ده۔ د
تعريف له مخې دا په دې معنا ده چې تابعه~$\Char{R}(x_0, \dots, x_k)$
محاسبه کېدونکې ده۔ د~\olref[int]{thm:representable-iff-comp} له
مخې، $\Char{R}$ په کوم فارمول، مثلاً
$!A_{\Char{R}}(x_0, \dots, x_k, y)$، تمثيلېږي۔
$!A_R(x_0, \dots, x_k)$ هغه فارمول ونيسئ چې
$!A_{\Char{R}}(x_0, \dots, x_k, \num{1})$ وي۔ بيا د هر~$n_0$،
\dots،~$n_k$ دپاره، که~$R(n_0, \dots, n_k)$ رښتيا وي، نو
$\Char{R}(n_0, \dots, n_k) = 1$؛ په دې حالت کښې~$\Th{Q}$،
$!A_{\Char{R}}(\num{n_0}, \dots, \num{n_k}, \num{1})$ ثابتوي، او
نو~$\Th{Q}$، $!A_R(\num{n_0}, \dots, \num{n_k})$ ثابتوي۔ له بلې
خوا، که~$R(n_0, \dots, n_k)$ ناسم وي، نو
$\Char{R}(n_0, \dots, n_k) = 0$۔ دا په دې معنا ده چې~$\Th{Q}$
لاندې ثابتوي:
\[
\lforall[y][(!A_{\Char{R}}(\num{n_0}, \dots, \num{n_k}, y) \lif y =
  \num{0})].
\]
څرنګه چې~$\Th{Q}$، $\eq/[\num{0}][\num{1}]$ ثابتوي، $\Th{Q}$،
$\lnot !A_{\Char{R}}(\num{n_0}, \dots, \num{n_k}, \num{1})$
ثابتوي، او نو~$\lnot !A_R(\num{n_0}, \dots, \num{n_k})$ هم
ثابتوي۔
\end{proof}

\begin{prob}
وښيئ چې که~$R$ په~$\Th{Q}$ کښې تمثيلېدونکې وي، نو~$\Char{R}$ هم
همداسې ده۔
\end{prob}

\end{document}
